% Luc Destombe --- licence CC BY-NC-SA 4.0
% https://creativecommons.org/licenses/by-nc-sa/4.0/deed.fr
% Fichier autonome : compiler avec pdflatex, deux fois (signets, nombre de pages).
\documentclass[10pt,a4paper,twocolumn]{article}
\makeatletter
\newif\iflycee@unecolonne
\newif\iflycee@palatino
\lycee@palatinotrue

\RequirePackage[utf8]{inputenc}
\RequirePackage[T1]{fontenc}
\RequirePackage[french]{babel}
\RequirePackage{etoolbox}

\newcommand{\ShorthandsOff}{\@ifundefined{shorthandoff}{}{\shorthandoff{;:!?}}}
\newcommand{\ShorthandsOn}{\@ifundefined{shorthandon}{}{\shorthandon{;:!?}}}
\ShorthandsOff
\AtBeginDocument{\ShorthandsOn}
\AtBeginEnvironment{tikzpicture}{\ShorthandsOff}

\RequirePackage{geometry}
\RequirePackage{fancyhdr}
\RequirePackage{lastpage}
\RequirePackage{multicol}
\RequirePackage{array}
\RequirePackage{enumitem}
\RequirePackage{xcolor}
\RequirePackage{needspace}

\RequirePackage{amsmath,amssymb}
\RequirePackage{mathpazo}
\DeclareMathSizes{10.95}{10.95}{8}{5.5}
\begingroup\catcode`\_=\active \catcode`\^=\active
\gdef_#1{\sb{\scriptscriptstyle #1}}
\gdef^#1{\sp{\scriptscriptstyle\lycee@moinsepais #1}}
\endgroup
\newcommand\lycee@moins{\mathbin{%
  \setbox\z@\hbox{$\m@th\scriptscriptstyle\lycee@moinsorig$}%
  \dimen@=.55\wd\z@ \advance\dimen@-.5pt
  \hbox to.75\wd\z@{\hss$\m@th\scriptscriptstyle\vcenter{\hbox{%
    \tikz\draw[line width=.5pt,line cap=round](0,0)--(\the\dimen@,0);}}$\hss}}}
\begingroup\lccode`\~=`\- \lowercase{\endgroup\let~}\lycee@moins
\newcommand\lycee@moinsepais{\mathcode`\-="8000 }
\AtBeginDocument{\mathchardef\lycee@moinsorig=\mathcode`\-\relax}
\AtBeginDocument{\catcode`\_=12 \mathcode`\_="8000
  \catcode`\^=12 \mathcode`\^="8000 }
\newcommand\lycee@exposants{%
  \fontdimen13\textfont2=.48\fontdimen6\textfont2
  \fontdimen14\textfont2=.48\fontdimen6\textfont2
  \fontdimen15\textfont2=.48\fontdimen6\textfont2
  \fontdimen13\scriptfont2=.48\fontdimen6\scriptfont2
  \fontdimen14\scriptfont2=.48\fontdimen6\scriptfont2
  \fontdimen15\scriptfont2=.48\fontdimen6\scriptfont2}
\every@math@size\expandafter{\the\every@math@size\lycee@exposants}
\newlength{\retraitcalculs}
\setlength{\retraitcalculs}{2em}

\RequirePackage{tikz}
\usetikzlibrary{arrows.meta,calc,babel,shapes.misc}
\RequirePackage[most]{tcolorbox}
\tcbuselibrary{skins,breakable}

\definecolor{coulDef}{HTML}{1F4E79}
\definecolor{coulProp}{HTML}{7030A0}
\definecolor{coulRem}{HTML}{C55A11}
\definecolor{coulEx}{HTML}{2E7D32}
\definecolor{coulExo}{HTML}{B03A2E}
\definecolor{coulPart}{HTML}{1F4E79}
\definecolor{coulMeth}{HTML}{1B6E4F}
\definecolor{coulCourbe}{HTML}{0B5ED7}
\definecolor{coulCourbeB}{HTML}{C00000}
\definecolor{coulCourbeC}{HTML}{2E7D32}
\definecolor{coulGrille}{HTML}{8C8C8C}
\definecolor{coulRep}{HTML}{1B6E4F}
\definecolor{coulBar}{HTML}{2E7D32}

\newcommand{\R}{\mathbb{R}}
\newcommand{\Rs}{\mathbb{R}^{*}}
\renewcommand{\le}{\leqslant}
\renewcommand{\ge}{\geqslant}
\renewcommand{\approx}{\simeq}
\newcommand{\vect}[1]{\overrightarrow{#1}}
\newcommand{\norme}[1]{\left\lVert #1 \right\rVert}
\newcommand{\intervalle}[2]{\left[#1\,;#2\right]}
\RequirePackage{cancel}
\renewcommand{\CancelColor}{\color{coulExo}}
\newcommand{\fois}[1]{\times{\color{coulExo}#1}}
\newcommand{\barre}[1]{\cancel{#1}}

\tikzset{grille/.style={coulGrille,line width=0.4pt}}
\newcommand{\quadrillage}[4]{\draw[grille,step=1] (#1,#2) grid (#3,#4);}
\tikzset{
  croix/.style={cross out,draw,line width=0.6pt,inner sep=0pt,minimum size=4pt},
  croixdroite/.style={croix,rotate=45,line width=0.8pt},
}
\newcommand{\pt}[3]{\path (#1) node[croix]{} node[#2,font=\small,inner sep=2.5pt]{$#3$};}

\newcommand{\lycee@repere}[4]{%
  \draw[grille,step=1] (#1,#2) grid (#3,#4);
  \draw[-{Stealth[length=2mm]},thick] (#1,0)--({#3+0.4},0) node[below left=-1pt]{$x$};
  \draw[-{Stealth[length=2mm]},thick] (0,#2)--(0,{#4+0.4}) node[below left=-1pt]{$y$};
  \node[below left,font=\scriptsize,inner sep=1.5pt] at (0,0) {$0$};}
\newcommand{\lycee@gradx}[1]{\foreach \k/\l in {#1}{%
  \draw (\k,0.07)--(\k,-0.07) node[below,font=\scriptsize,inner sep=1.5pt]{$\l$};}}
\newcommand{\lycee@grady}[1]{\foreach \k/\l in {#1}{%
  \draw (0.07,\k)--(-0.07,\k) node[left,font=\scriptsize,inner sep=1.5pt]{$\l$};}}
\newcommand{\lycee@intff}[2]{\left[#1\,;#2\right]}
\newcommand{\lycee@intfo}[2]{\left[#1\,;#2\right[}
\newcommand{\lycee@intof}[2]{\left]#1\,;#2\right]}
\newcommand{\lycee@intoo}[2]{\left]#1\,;#2\right[}
\newcommand{\lycee@ens}[1]{\left\{#1\right\}}
\AtBeginDocument{%
  \providecommand{\repere}{\lycee@repere}%
  \providecommand{\gradx}{\lycee@gradx}%
  \providecommand{\grady}{\lycee@grady}%
  \providecommand{\Cf}{\mathcal{C}\sb{\scriptscriptstyle f}}%
  \providecommand{\Cg}{\mathcal{C}\sb{\scriptscriptstyle g}}%
  \providecommand{\intff}{\lycee@intff}%
  \providecommand{\intfo}{\lycee@intfo}%
  \providecommand{\intof}{\lycee@intof}%
  \providecommand{\intoo}{\lycee@intoo}%
  \providecommand{\ens}{\lycee@ens}}

\newlist{questions}{enumerate}{3}
\setlist[questions,1]{label=\arabic*\textdegree),leftmargin=*,itemsep=2pt,topsep=3pt}
\setlist[questions,2]{label=\alph*),leftmargin=*,itemsep=1pt,topsep=2pt}
\setlist[questions,3]{label=\roman*),leftmargin=*,itemsep=1pt,topsep=1pt}
\setlist[itemize]{label=\textbullet,leftmargin=*,itemsep=2pt,topsep=3pt}

\newcommand{\besoinplace}[1]{\par\penalty\@M\begingroup
  \dimen@=#1\relax\dimen@ii=\pagegoal\advance\dimen@ii-\pagetotal
  \ifdim\dimen@>\dimen@ii\ifdim\dimen@ii>\z@\vfil\fi\break\fi\endgroup}

\newcommand{\lycee@pied}{\thepage/\pageref{LastPage}}
\newcommand{\entete}[2]{%
  \pagestyle{fancy}\fancyhf{}%
  \fancyhead[L]{\footnotesize\color{coulPart}\textbf{#1}}%
  \fancyhead[R]{\footnotesize\color{coulPart}\textbf{#2}}%
  \fancyfoot[C]{\footnotesize\lycee@pied}%
  \renewcommand{\headrulewidth}{0.4pt}%
  \renewcommand{\headrule}{\hbox to\headwidth{%
    \color{coulPart!40}\leaders\hrule height \headrulewidth\hfill}}}

\RequirePackage{ccicons}
\newcommand{\lycee@licence}{{\scriptsize\color{gray}%
  \copyright\ Luc Destombe --- \ccbyncsa\ CC BY-NC-SA 4.0}}
\newif\iflycee@licence \lycee@licencetrue
\newcommand{\sanslicence}{\lycee@licencefalse}
\AtBeginDocument{\iflycee@licence\AddToHookNext{shipout/foreground}{%
  \put(0,0){\rlap{\kern\dimexpr1in+\hoffset+\oddsidemargin+\textwidth\relax
    \llap{\raisebox{-\dimexpr1in+\voffset+\topmargin+\headheight+\headsep
      +\textheight+\footskip\relax}{\lycee@licence}}}}}\fi}

\newcommand{\formulesserrees}{%
  \setlength{\abovedisplayskip}{3pt}\setlength{\belowdisplayskip}{3pt}%
  \setlength{\abovedisplayshortskip}{2pt}\setlength{\belowdisplayshortskip}{3pt}}

\setlength{\parindent}{0pt}

\RequirePackage[implicit=false,hidelinks,pdfencoding=auto,psdextra,bookmarksopen,
  bookmarksopenlevel=1,pdfcreator={},pdfproducer={}]{hyperref}
\RequirePackage{bookmark}
\hypersetup{pdfauthor={Luc Destombe},pdfkeywords={CC BY-NC-SA 4.0}}
\newcounter{lycee@signet}
\newcommand{\signet}[3][]{\stepcounter{lycee@signet}%
  \edef\lycee@dest{\if\relax\detokenize{#1}\relax
    signet.\arabic{lycee@signet}\else#1\fi}%
  \hypertarget{\lycee@dest}{}\bookmark[level=#2,dest=\lycee@dest]{#3}}
\pdfstringdefDisableCommands{%
  \def\R{R}\def\vect#1{#1}\def\Cf{Cf}\def\Cg{Cg}\def\fois#1{×#1}%
  \def\barre#1{#1}\def\intervalle#1#2{[#1;#2]}\def\intff#1#2{[#1;#2]}%
  \def\textsuperscript#1{#1}\def\dfrac#1#2{#1/#2}\def\frac#1#2{#1/#2}%
  \def\vec#1{#1⃗}}%
\RequirePackage[official]{eurosym}

\tcbset{exercice corrige/.style={
  enhanced,breakable,before upper=\raggedright,
  colback=white,colframe=coulExo,
  boxrule=0.8pt,arc=2pt,
  left=6pt,right=6pt,top=8pt,bottom=5pt,
  before skip=9pt,after skip=4pt,
  coltitle=white,fonttitle=\bfseries\small,
  attach boxed title to top left={xshift=8pt,yshift=-\tcboxedtitleheight/2},
  boxed title style={colback=coulExo,arc=2pt,boxrule=0pt}}}
\newcounter{exo}
\newtcolorbox{exercice}[1][]{exercice corrige,
  phantom={\signet[exercice-\arabic{exo}]{2}{Exercice \arabic{exo}}},
  title={Exercice \theexo\ifx\relax#1\relax\else{} --- #1\fi}}
\newenvironment{exo}[1][\relax]{\refstepcounter{exo}\begin{exercice}[#1]}{\end{exercice}}

\RequirePackage{cuted}
\geometry{top=1.6cm,bottom=1cm,left=1cm,right=1cm,headheight=14pt,footskip=0.6cm}

\newcommand{\entetecorrige}[3]{%
  \begin{strip}
  \begin{center}
    \begin{tcolorbox}[enhanced,width=0.92\linewidth,colback=white,colframe=coulPart,
        boxrule=1.6pt,arc=3pt,halign=center,top=5pt,bottom=5pt]
      {\LARGE\bfseries\color{coulPart} #1}\\[3pt]
      {\normalsize #2}
    \end{tcolorbox}
  \end{center}
  \if\relax\detokenize{#3}\relax\else

  \vspace{2pt}
  \noindent
  \begin{tcolorbox}[enhanced,colback=white,colframe=black!35,boxrule=0.5pt,arc=2pt,
      left=7pt,right=7pt,top=4pt,bottom=4pt]
  {\small #3}
  \end{tcolorbox}
  \vspace{6pt}
  \fi
  \end{strip}}

\newcounter{partie}
\newcommand{\partiebandeau}[1]{%
  \besoinplace{8\baselineskip}\vspace{6pt}%
  \begin{tcolorbox}[enhanced,colback=coulPart!8,colframe=coulPart,boxrule=0pt,
      leftrule=3pt,arc=0pt,left=5pt,right=4pt,top=3pt,bottom=3pt,
      before skip=4pt,after skip=2pt]
    \bfseries\color{coulPart}#1\signet{1}{#1}
  \end{tcolorbox}}
\newcommand{\partie}[1]{\stepcounter{partie}\partiebandeau{\Roman{partie}) #1}}
\newcommand{\partiesansnum}[1]{\partiebandeau{#1}}

\newcommand{\res}[1]{%
  \tcbox[on line,colback=coulPart!7,colframe=coulPart,boxrule=0.5pt,arc=1pt,
    boxsep=0pt,left=3pt,right=3pt,top=1pt,bottom=2pt]{$#1$}}
\newcommand{\calc}[1]{\par\smallskip\noindent$\displaystyle #1$\par}
\newenvironment{calculs}{\par\smallskip\noindent$\displaystyle\begin{aligned}[t]}%
  {\end{aligned}$\par\smallskip}
\newcommand{\rem}[1]{\par\smallskip{\small\itshape\color{coulPart}#1}\par}
\newcommand{\deux}[2]{\par\noindent\makebox[0.5\linewidth][l]{#1}#2}

\setlist[questions,1]{label=\arabic*\textdegree),leftmargin=*,itemsep=2pt,topsep=2pt}
\setlist[questions,2]{label=\alph*),leftmargin=*,itemsep=2pt,topsep=1pt}
\newlist{reponses}{enumerate}{1}
\setlist[reponses]{label=\alph*),leftmargin=*,itemsep=2pt,topsep=2pt}

\setlength{\parskip}{1pt}
\setlength{\columnsep}{0.6cm}
\raggedbottom
\formulesserrees
\AtBeginDocument{\formulesserrees}

\makeatother
\entete{Chapitre 2 --- Fonctions --- Corrigé des exercices}{1\textsuperscript{re} STI2D}

\usepackage{tkz-tab}
\definecolor{coulCourbe}{HTML}{C0392B}
\definecolor{coulCourbeB}{HTML}{1F4E79}
\let\deuxpaquet\deux
\renewcommand{\deux}[2]{\deuxpaquet{#1}{#2}\par}
\newenvironment{tabsignes}[2]{\begin{center}\begin{tikzpicture}[scale=0.58]
  \tkzTabInit[lgt=2.2,espcl=1.7]{#1}{#2}}{\end{tikzpicture}\end{center}}
\clubpenalty=10000
\widowpenalty=10000

\begin{document}

\entetecorrige{CORRIGÉ --- FONCTIONS}
  {Chapitre 2 : fiche d'exercices --- Première STI2D}
  {Les résultats sont encadrés ; les remarques en bleu signalent les erreurs
   fréquentes et les méthodes à retenir. Quand on demande « une courbe qui peut
   représenter $f$ », la courbe tracée n'est qu'une réponse possible : seuls comptent
   les points du tableau et le sens de variation.}

\partie{Antécédent et image}

\begin{exo}[Images et antécédent]
\begin{questions}
  \item \begin{calculs}
        f(1)&=4-2\\
        f(1)&=\res{2}\\
        f(-2)&=-8-2\times 4\\
        f(-2)&=\res{-16}\\
        f(2{,}5)&=10-2\times 6{,}25\\
        f(2{,}5)&=\res{-2{,}5}
        \end{calculs}
  \item $f(3)=12-18$, soit $-6$ : \res{\text{oui}}, $3$ est un antécédent de $-6$.
\end{questions}
\rem{Dans $-2x^{2}$ pour $x=-2$ : d'abord le carré, $(-2)^{2}=4$, puis
$-2\times 4=-8$.}
\end{exo}

\begin{exo}[Une fonction affine]
\begin{questions}
  \item $g(-3)=-15-8=\res{-23}$ ; \ $g(4)=20-8=\res{12}$ ;\\
        $g\big(\frac{6}{5}\big)=6-8=\res{-2}$
  \item $g\big(\frac{2}{5}\big)=2-8=-6$, et $-6\ne 6$ : \res{\text{non}}.
\end{questions}
\end{exo}

\begin{exo}[Plusieurs antécédents]
\begin{questions}
  \item $h(0)=\res{-6}$ ; \ $h(2)=4-2-6=\res{-4}$ ; \ $h(-2)=4+2-6=\res{0}$
  \item $h(3)=9-3-6=0$ : \res{\text{oui}}.
  \item \res{-2\text{ et }3}, d'après les questions 1 et 2.
\end{questions}
\end{exo}

\begin{exo}[Calculer des antécédents]
\rem{Chercher les antécédents de $y$, c'est résoudre l'équation $f(x)=y$.}
\begin{reponses}
  \item $x-4=3$, soit $x=7$ : \ \res{7}
  \item $x^{2}=36$ : \ \res{-6\text{ et }6}
  \item $\frac{1}{x}=4$, soit $1=4x$ : \ \res{\tfrac{1}{4}}
  \item $x(x-5)=0$ (produit nul) : \ \res{0\text{ et }5}
\end{reponses}
\rem{En b), ne pas oublier $-6$ : $(-6)^{2}=36$ aussi.}
\end{exo}

\partie{Ensemble de définition}

\begin{exo}[Valeurs interdites]
\begin{itemize}
  \item $f$ : \res{0}, on ne divise pas par $0$.
  \item $g$ : \res{5}, le dénominateur $x-5$ s'annule.
  \item $h$ : \res{-2}, la racine d'un nombre négatif n'existe pas.
\end{itemize}
\end{exo}

\begin{exo}[Calculer, si c'est possible]
\begin{reponses}
  \item $f(6)=\frac{8}{2}=\res{4}$ ; \ $f(-4)=\frac{8}{-8}=\res{-1}$ ; \ $f(4)=\frac{8}{0}$
        \res{\text{n'existe pas}}.
  \item $g(7)=\sqrt{9}=\res{3}$ ; \ $g(14)=\sqrt{16}=\res{4}$ ; \ $g(-3)=\sqrt{-1}$
        \res{\text{n'existe pas}}.
\end{reponses}
\end{exo}

\begin{exo}[Fractions]
\rem{On cherche la valeur qui annule le dénominateur.}
\deux{a) \res{\R\setminus\{-3\}}}{b) \res{\R\setminus\{9\}}}
\deux{c) $3x=12$ : \res{\R\setminus\{4\}}}{d) $4x=-20$ : \res{\R\setminus\{-5\}}}
\end{exo}

\begin{exo}[Racines carrées]
\rem{Sous la racine : un nombre positif ou nul.}
\deux{a) $x\ge 4$ : \res{\left[4\,;+\infty\right[}}{b) $x\ge -6$ : \res{\left[-6\,;+\infty\right[}}
\deux{c) $2x\ge 10$ : \res{\left[5\,;+\infty\right[}}{d) $3x\ge -5$ : \res{\left[-\tfrac{5}{3}\,;+\infty\right[}}
\end{exo}

\partie{Tableau de valeurs, courbe et lecture graphique}

\begin{exo}[Première lecture graphique]
\begin{reponses}
  \item $f(0)=\res{3}$ et $f(2)=\res{3}$
  \item L'image de $-1$ est \res{0}.
  \item $4$ a pour antécédent \res{1} (le point le plus haut).
  \item Les antécédents de $0$ sont \res{-1} et \res{3}.
  \item \res{\text{Non}} : la courbe ne monte pas au-dessus de $4$ ; la droite $y=5$ ne
        la coupe pas.
\end{reponses}
\rem{Image : on part de l'axe des abscisses. Antécédents : on part de l'axe des
ordonnées et on lit \emph{toutes} les abscisses des points d'intersection.}
\end{exo}

\begin{exo}[Lire un tableau de valeurs]
\begin{questions}
  \item $f(-1)=\res{2}$ ; \ $f(1)=\res{-3}$ ; \ $f(-2)=\res{4}$
  \item $f(1)=-3$ : \res{\text{oui}}.
  \item $f(2)=1$ : \res{2} est un antécédent de $1$.
\end{questions}
\rem{Ligne du haut : les nombres $x$ ; ligne du bas : leurs images.}
\end{exo}

\begin{exo}[Images et antécédents dans un tableau]
\begin{questions}
  \item $g(4)=\res{-1}$ ; \ $g(2)=\res{0}$ ; \ $g(-1)=\res{6}$
  \item $6$ : \res{-1} ; \ $-1$ : \res{4} ; \ $-4$ : \res{-3} ou \res{6}
  \item $-4$ a deux antécédents dans ce tableau ($-3$ et $6$) : on ne peut pas parler
        de « l'» antécédent.
\end{questions}
\rem{Un nombre a une seule image, mais il peut avoir plusieurs antécédents.}
\end{exo}

\begin{exo}[Lecture graphique complète]
\begin{multicols}{2}
\begin{reponses}
  \item \res{\intff{-4}{5}}
  \item $f(3)=\res{-1}$
  \item $f(-1)=\res{1}$
  \item $f(1)=\res{1}$ ; $f(-4)=\res{-2}$
  \item \res{0}
  \item \res{-1\,;1\,;5}
  \item \res{-2\,;2\,;4}
  \item \res{-3\,;3}
  \item \res{\text{aucun}}
\end{reponses}
\end{multicols}
\begin{reponses}[start=10]
  \item Maximum \res{2}, atteint en $x=0$ ; minimum \res{-2}, atteint en $x=-4$.
\end{reponses}
\rem{Pour $1$ : la droite $y=1$ coupe la courbe en trois points ; ne pas oublier celui
d'abscisse $5$, au bout de la courbe. Le minimum est atteint à une borne.}
\end{exo}

\begin{exo}[Deux expressions, une seule fonction]
\begin{questions}
  \item $f(2)=6\times 1+4-5=\res{5}$ \ et \ $g(2)=\res{5}$
  \item Par exemple pour $x=0$ : $f(0)=12-5=7$ et $g(0)=7$.
  \item On développe :
        \begin{calculs}
        f(x)&=3x-x^{2}+12-4x+x^{2}-5\\
        f(x)&=-x+7
        \end{calculs}
        donc \res{f(x)=g(x)} pour tout nombre $x$.
\end{questions}
\rem{Des exemples (questions 1 et 2) ne démontrent rien : seul le calcul de la
question 3 vaut pour \emph{tous} les nombres.}
\end{exo}

\partie{Tableau de variations d'une fonction}

\begin{exo}[De la courbe au tableau]
$f$ est décroissante sur $\intff{-3}{1}$, puis croissante sur $\intff{1}{5}$.
\begin{center}
\begin{tikzpicture}[scale=0.62]
  \tkzTabInit[lgt=1.4,espcl=2.4]{$x$/1,$f$/2}{$-3$,$1$,$5$}
  \tkzTabVar{+/$3$,-/$-2$,+/$4$}
\end{tikzpicture}
\end{center}
Maximum \res{4}, atteint en $x=5$ ; minimum \res{-2}, atteint en $x=1$.
\rem{Le maximum est la plus grande des deux valeurs aux bornes : $f(5)=4>f(-3)=3$.}
\end{exo}

\begin{exo}[Décrire puis résumer]
\begin{questions}
  \item $f$ est croissante sur $\intff{-3}{-1}$, décroissante sur $\intff{-1}{2}$, puis
        croissante sur $\intff{2}{4}$.
  \item \leavevmode
\begin{center}
\begin{tikzpicture}[scale=0.62]
  \tkzTabInit[lgt=1.4,espcl=1.8]{$x$/1,$f$/2}{$-3$,$-1$,$2$,$4$}
  \tkzTabVar{-/$-1$,+/$2$,-/$-3$,+/$0$}
\end{tikzpicture}
\end{center}
\end{questions}
\end{exo}

\begin{exo}[Du tableau à la courbe]
\begin{minipage}[c]{0.46\linewidth}\raggedright
\begin{questions}
  \item $f$ est définie sur \res{\intff{-4}{3}}.
  \item Une courbe possible ci-contre : elle part de $(-4\,;-2)$, monte jusqu'à
        $(0\,;3)$, puis descend jusqu'à $(3\,;-1)$.
\end{questions}
\end{minipage}\hfill
\begin{minipage}[c]{0.52\linewidth}\centering
\begin{tikzpicture}[x=0.3cm,y=0.3cm]
  \repere{-5}{-3}{4}{4}
  \gradx{-4,-2,2} \grady{-2,1,3}
  \draw[coulCourbe,line width=1pt] (-4,-2) .. controls (-2.5,1) and (-1.2,3) .. (0,3)
    .. controls (1.2,3) and (2.4,1) .. (3,-1);
  \foreach \p in {(-4,-2),(0,3),(3,-1)} \path[coulCourbe] \p node[croix]{};
\end{tikzpicture}
\end{minipage}
\end{exo}

\begin{exo}[Lire un tableau de variations]
\begin{questions}
  \item \res{\intff{-5}{4}}
  \item $g$ est décroissante sur $\intff{-5}{-2}$, croissante sur $\intff{-2}{1}$, puis
        décroissante sur $\intff{1}{4}$.
  \item Maximum \res{4}, atteint en $x=1$ ; minimum \res{-2}, atteint en $x=-2$.
  \item Une courbe possible :
\end{questions}
\begin{center}
\begin{tikzpicture}[x=0.3cm,y=0.3cm]
  \repere{-6}{-3}{5}{5}
  \gradx{-5,-2,1,4} \grady{-2,3,4}
  \draw[coulCourbe,line width=1pt] (-5,3) .. controls (-4,1) and (-3,-2) .. (-2,-2)
    .. controls (-1,-2) and (0,4) .. (1,4) .. controls (2,4) and (3,1.5) .. (4,0);
  \foreach \p in {(-5,3),(-2,-2),(1,4),(4,0)} \path[coulCourbe] \p node[croix]{};
\end{tikzpicture}
\end{center}
\rem{On compare les valeurs aux bornes aux extremums intérieurs : $g(-5)=3<4$ et
$g(4)=0>-2$.}
\end{exo}

\begin{exo}[Construire un tableau à partir d'un texte]
\begin{questions}
  \item \leavevmode
\begin{center}
\begin{tikzpicture}[scale=0.62]
  \tkzTabInit[lgt=1.4,espcl=1.8]{$x$/1,$f$/2}{$-4$,$1$,$3$,$5$}
  \tkzTabVar{-/$-3$,+/$2$,-/$-1$,+/$2$}
\end{tikzpicture}
\end{center}
  \item Une courbe possible :
\end{questions}
\begin{center}
\begin{tikzpicture}[x=0.3cm,y=0.3cm]
  \repere{-5}{-4}{6}{3}
  \gradx{-4,1,3,5} \grady{-3,-1,2}
  \draw[coulCourbe,line width=1pt] (-4,-3) .. controls (-2,-1) and (-0.5,2) .. (1,2)
    .. controls (2,2) and (2.3,-1) .. (3,-1) .. controls (3.8,-1) and (4.5,1) .. (5,2);
  \foreach \p in {(-4,-3),(1,2),(3,-1),(5,2)} \path[coulCourbe] \p node[croix]{};
\end{tikzpicture}
\end{center}
\end{exo}

\begin{exo}[Résolution graphique]
\rem{$f(x)=k$ : abscisses des points de la courbe situés \emph{sur} la droite $y=k$ ;
$f(x)>k$ : abscisses des points situés \emph{au-dessus}.}
\begin{questions}
  \item \begin{questions}
          \item \res{S=\ens{0\,;8\,;10}}
          \item \res{S=\ens{1\,;7}}
          \item \res{S=\ens{4}}
          \item \res{S=\varnothing} : le maximum de $f$ est $3$.
          \item \res{S=\ens{2\,;6}}
        \end{questions}
  \item La droite $y=-1$ coupe la courbe en deux points, d'abscisses $3$ et $5$ :
        \res{\text{deux solutions}}.
  \item \begin{questions}
          \item \res{S=\left[0\,;4\right[\cup\left]4\,;10\right]}
          \item \res{S=\left]1\,;7\right[}
          \item \res{S=\intff{0}{1}\cup\intff{7}{10}}
          \item \res{S=\intff{0}{2}\cup\intff{6}{10}}
        \end{questions}
\end{questions}
\rem{En 3) a), seul $4$ est exclu : $f(4)=-2$ n'est pas strictement supérieur à $-2$.}
\end{exo}

\partie{Fonctions affines --- équation réduite}

\begin{exo}[Reconnaître $m$ et $p$]
\deux{a) $m=\res{5}$ ; $p=\res{-2}$}{b) $m=\res{-3}$ ; $p=\res{6}$}
\deux{c) $m=\res{1}$ ; $p=\res{-4}$}{d) $m=\res{-3}$ ; $p=\res{7}$}
\deux{e) $m=\res{0}$ ; $p=\res{4}$}{f) $m=\res{-2}$ ; $p=\res{0}$}
\deux{g) $m=\res{-\tfrac{1}{3}}$ ; $p=\res{1}$}{h) $m=\res{\tfrac{2}{5}}$ ; $p=\res{-\tfrac{3}{5}}$}
\rem{En d), $m$ est le nombre devant $x$, même écrit en second. En g),
$\frac{-x}{3}=-\frac{1}{3}x$.}
\end{exo}

\begin{exo}[Images et ordonnée à l'origine]
\begin{questions}
  \item $f(4)=-8+5=\res{-3}$ \ et \ $f(-3)=6+5=\res{11}$
  \item $p=f(0)=\res{5}$ : la droite coupe l'axe des ordonnées en $(0\,;5)$.
\end{questions}
\end{exo}

\begin{exo}[Pente et équation réduite (1)]
\begin{calculs}
m&=\frac{3-(-2)}{0-5}\\
m&=\res{-1}
\end{calculs}
$B(0\,;3)$ est sur l'axe des ordonnées : $p=3$. \ \res{y=-x+3}\par
Contrôle avec $A$ : $-5+3=-2$.
\end{exo}

\begin{exo}[Pente et équation réduite (2)]
\begin{calculs}
m&=\frac{10-1}{4-(-2)}\\
m&=\res{\tfrac{3}{2}}
\end{calculs}
Avec $A$ : $1=\frac{3}{2}\times(-2)+p$, soit $1=-3+p$ et $p=4$. \ \res{y=\tfrac{3}{2}x+4}\par
Contrôle avec $B$ : $6+4=10$.
\end{exo}

\partie{Tracer la droite d'équation $y=mx+p$}

\begin{exo}[Sens de variation]
\rem{Une fonction affine est croissante si $m>0$, décroissante si $m<0$.}
\begin{questions}
  \item $m=3>0$ : \res{\text{croissante}}
  \item $m=-2<0$ : \res{\text{décroissante}}
  \item $m=1>0$ : \res{\text{croissante}}
  \item $m=-1<0$ : \res{\text{décroissante}}
  \item $f(x)=\sqrt{2}\,x-3\sqrt{2}$, $m=\sqrt{2}>0$ : \res{\text{croissante}}
  \item $f(x)=-\frac{3}{5}x+\frac{4}{5}$, $m=-\frac{3}{5}<0$ : \res{\text{décroissante}}
\end{questions}
\end{exo}

\begin{exo}[Associer chaque fonction à sa droite]
\begin{itemize}
  \item \textcircled{1} monte et coupe l'axe des ordonnées en $-3$ : \res{g}, $m=2$,
        $p=-3$.
  \item \textcircled{2} descend doucement et coupe l'axe en $1$ : \res{h},
        $m=-\frac{1}{2}$, $p=1$.
  \item \textcircled{3} descend fortement et coupe l'axe en $2$ : \res{f}, $m=-3$,
        $p=2$.
\end{itemize}
\end{exo}

\begin{exo}[Appartenance et tracé]
\begin{minipage}[t]{0.58\linewidth}\raggedright
\begin{questions}
  \item $4\times 1-3=1$, c'est $y_{A}$ : \res{A\in d}.
  \item $4x-3=5$, soit $x=2$ : \res{B(2\,;5)}.
  \item On trace la droite qui passe par $(0\,;-3)$, $A$ et $B$.
\end{questions}
\end{minipage}\hfill
\begin{minipage}[t]{0.4\linewidth}\centering\vspace{0pt}
\begin{tikzpicture}[x=0.3cm,y=0.3cm]
  \repere{-2}{-4}{3}{5}
  \gradx{-1,1,2} \grady{-3,1,5}
  \draw[coulCourbeB,line width=1pt,domain=-0.25:2] plot (\x,{4*\x-3});
  \path[coulCourbe] (0,-3) node[croix]{} (1,1) node[croix]{} node[right,font=\tiny] {$A$}
    (2,5) node[croix]{} node[right,font=\tiny] {$B$};
\end{tikzpicture}
\end{minipage}
\end{exo}

\begin{exo}[Tracer à partir d'un point et de la pente]
\rem{Depuis $A$, on avance de $1$ et on monte de $m$ (on descend si $m<0$).}
\begin{minipage}[t]{0.58\linewidth}\raggedright
\begin{reponses}
  \item $2=3\times 1+p$, donc $p=-1$ :\\ \res{y=3x-1}
  \item $-1=-2\times 3+p$, donc $p=5$ :\\ \res{y=-2x+5}
\end{reponses}
\end{minipage}\hfill
\begin{minipage}[t]{0.4\linewidth}\centering\vspace{0pt}
\begin{tikzpicture}[x=0.3cm,y=0.3cm]
  \repere{-5}{-5}{5}{5}
  \gradx{-4,-2,2,4} \grady{-4,-2,2,4}
  \draw[coulCourbe,line width=1pt,domain=-4/3:2] plot (\x,{3*\x-1});
  \draw[coulCourbeB,line width=1pt,domain=0:5] plot (\x,{-2*\x+5});
  \path[coulCourbe] (1,2) node[croix]{}; \path[coulCourbeB] (3,-1) node[croix]{};
  \node[coulCourbe,font=\footnotesize,left] at (1.6,4.4) {a)};
  \node[coulCourbeB,font=\footnotesize,right] at (4,-3.2) {b)};
\end{tikzpicture}
\end{minipage}
\end{exo}

\partie{Équation cartésienne d'une droite}

\begin{exo}[Équation cartésienne de $(AB)$ (1)]
$\vec{AB}\,(5-2\,;1-3)$, soit $\vec{AB}\,(3\,;-2)$. Avec $\vec{u}\,(-b\,;a)$ : $b=-3$ et
$a=-2$, d'où $-2x-3y+c=0$.\par
Avec $A$ : $-4-9+c=0$, $c=13$. \ \res{-2x-3y+13=0} \ ou \ \res{2x+3y-13=0}\par
Contrôle : $4+9-13=0$ et $10+3-13=0$.
\rem{Dans $\vec{u}\,(-b\,;a)$, la première coordonnée donne $-b$ : attention au signe
de $b$.}
\end{exo}

\begin{exo}[Équation cartésienne de $(AB)$ (2)]
$\vec{AB}\,(8\,;4)$ : $b=-8$ et $a=4$, d'où $4x-8y+c=0$. Avec $A$ : $-12+8+c=0$, $c=4$.\par
\res{4x-8y+4=0}, soit, en divisant par $4$ : \res{x-2y+1=0}.\par
Contrôle : $-3+2+1=0$ et $5-6+1=0$.
\end{exo}

\begin{exo}[Lire une équation cartésienne]
\begin{questions}
  \item $A$ : $3-8+5=0$ : \res{A\in d}.\\
        $B$ : $-9+4+5=0$ : \res{B\in d}.\\
        $C$ : $6-12+5=-1\ne 0$ : \res{C\notin d}.
  \item $15-4y+5=0$, soit $4y=20$ : \res{y=5}.
  \item $3x+16+5=0$, soit $3x=-21$ : \res{x=-7}.
  \item $\vec{u}\,(-b\,;a)$ avec $a=3$ et $b=-4$ : \res{\vec{u}\,(4\,;3)}.
\end{questions}
\end{exo}

\partie{Tracer une droite d'équation $ax+by+c=0$}

\begin{exo}[Point et vecteur directeur]
\begin{minipage}[t]{0.58\linewidth}\raggedright
\begin{questions}
  \item $b=-1$ et $a=3$ : $3x-y+c=0$ ; avec $A$ : $3+2+c=0$, $c=-5$.\\
        \res{d : 3x-y-5=0}\\
        On trace par $A(1\,;-2)$ et $(2\,;1)$.
  \item $\vec{BC}\,(6\,;3)$ : $b=-6$ et $a=3$, d'où $3x-6y+c=0$ ; avec $B$ :
        $-12+18+c=0$, $c=-6$.\\
        \res{(BC) : x-2y-2=0}
\end{questions}
\end{minipage}\hfill
\begin{minipage}[t]{0.4\linewidth}\centering\vspace{0pt}
\begin{tikzpicture}[x=0.3cm,y=0.3cm]
  \repere{-5}{-5}{4}{4}
  \gradx{-4,1,2} \grady{-3,-2,1}
  \draw[coulCourbe,line width=1pt,domain=0:3] plot (\x,{3*\x-5});
  \draw[coulCourbeB,line width=1pt,domain=-5:4] plot (\x,{(\x-2)/2});
  \path[coulCourbe] (1,-2) node[croix]{} (2,1) node[croix]{};
  \path[coulCourbeB] (-4,-3) node[croix]{} (2,0) node[croix]{};
  \node[coulCourbe,font=\tiny] at (3.5,3.2) {$d$};
  \node[coulCourbeB,font=\tiny] at (-3.6,-1.6) {$(BC)$};
\end{tikzpicture}
\end{minipage}
\end{exo}

\begin{exo}[Deux droites dans un même repère]
\begin{minipage}[t]{0.56\linewidth}\raggedright
$d$ : pour $x=0$, $y=2$ ; pour $x=3$, $y=0$.\par
$d'$ : pour $x=0$, $y=-1$ ; pour $x=3$, $12-3y-3=0$, donc $y=3$.
\rem{Pour $d'$, on essaie des valeurs de $x$ jusqu'à ce que $4x-3$ soit un multiple
de $3$.}
\end{minipage}\hfill
\begin{minipage}[t]{0.42\linewidth}\centering\vspace{0pt}
\begin{tikzpicture}[x=0.28cm,y=0.28cm]
  \repere{-4}{-5}{5}{5}
  \gradx{-3,3} \grady{-1,2,3}
  \draw[coulCourbe,line width=1pt,domain=-3:5] plot (\x,{(6-2*\x)/3});
  \draw[coulCourbeB,line width=1pt,domain=-3:4.5] plot (\x,{(4*\x-3)/3});
  \foreach \p in {(0,2),(3,0)} \path[coulCourbe] \p node[croix]{};
  \foreach \p in {(0,-1),(3,3)} \path[coulCourbeB] \p node[croix]{};
  \node[coulCourbe,font=\tiny] at (-3.4,3.2) {$d$};
  \node[coulCourbeB,font=\tiny] at (4.2,3.4) {$d'$};
\end{tikzpicture}
\end{minipage}
\end{exo}

\begin{exo}[Droites verticales et horizontales]
\begin{minipage}[t]{0.52\linewidth}\raggedright
\begin{reponses}
  \item $x=-2$ : \res{\text{verticale}}
  \item $y=4$ : \res{\text{horizontale}}
  \item $y=x-1$ : $(0\,;-1)$ et $(1\,;0)$
  \item $(0\,;2)$ et $(2\,;-1)$
\end{reponses}
\end{minipage}\hfill
\begin{minipage}[t]{0.46\linewidth}\centering\vspace{0pt}
\begin{tikzpicture}[x=0.28cm,y=0.28cm]
  \repere{-4}{-4}{5}{5}
  \gradx{-2,2} \grady{-1,2,4}
  \draw[coulCourbe,line width=1pt] (-2,-4) -- (-2,5) node[below right,font=\tiny] {a};
  \draw[coulCourbeB,line width=1pt] (-4,4) -- (5,4) node[below left,font=\tiny] {b};
  \draw[coulCourbe!60!black,line width=1pt] (-3,-4) -- (5,4)
    node[below right,font=\tiny] {c};
  \draw[coulCourbeB!60!black,line width=1pt,domain=-4/3:4] plot (\x,{2-1.5*\x})
    node[right,font=\tiny] {d};
\end{tikzpicture}
\end{minipage}
\par\smallskip
Ont une équation réduite : \res{b,\ c,\ d} ($y=4$, $y=x-1$,
$y=-\frac{3}{2}x+2$) ; seule la verticale a) n'en a pas.
\end{exo}

\partie{Passer d'une forme à l'autre}

\begin{exo}[Pente, équation réduite, équation cartésienne (1)]
\begin{questions}
  \item $m=\dfrac{2-(-4)}{2-(-1)}=2$ ; avec $A$ : $-4=-2+p$, $p=-2$.\\
        \res{y=2x-2}
  \item On met tout dans un membre : \res{2x-y-2=0}.
\end{questions}
\end{exo}

\begin{exo}[Pente, équation réduite, équation cartésienne (2)]
$m=\dfrac{1-3}{4-0}=-\dfrac{1}{2}$ et $p=3$ ($A$ est sur l'axe des ordonnées).\par
\res{y=-\tfrac{1}{2}x+3} ; \ $\frac{1}{2}x+y-3=0$, puis, en multipliant par $2$ :
\res{x+2y-6=0}.
\end{exo}

\begin{exo}[Deux droites à déterminer]
\begin{reponses}
  \item $m=\dfrac{6-(-2)}{1-3}=-4$ ; avec $C$ : $-2=-12+p$, $p=10$ :
        \ \res{y=-4x+10}
  \item $5=5\times(-1)+p$, $p=10$ : \ \res{y=5x+10}
\end{reponses}
\end{exo}

\begin{exo}[Un point et une pente]
\begin{reponses}
  \item $-3=4\times 2+p$, $p=-11$ : \ \res{y=4x-11}, \ soit \ \res{4x-y-11=0}
  \item Pente nulle : droite horizontale qui passe par $B$ : \ \res{y=-4}, \ soit
        \ \res{y+4=0}
\end{reponses}
\end{exo}

\partie{Fonction carré}

\begin{exo}[Images et antécédents par la fonction carré]
\begin{questions}
  \item \res{16} ; \ \res{49} ; \ \res{0{,}04} ; \ \res{400}
  \item \res{-5\text{ et }5} ; \ \res{0} ; \ \res{\text{aucun}} ;
        \ \res{-\sqrt{11}\text{ et }\sqrt{11}}
\end{questions}
\rem{Un carré n'est jamais négatif : $-9$ n'a pas d'antécédent.}
\end{exo}

\begin{exo}[Lecture graphique avec la fonction carré]
\begin{questions}
  \item {\small\setlength{\tabcolsep}{2.5pt}\renewcommand{\arraystretch}{1.3}%
        \begin{tabular}{|c|c|c|c|c|c|c|c|c|c|}
        \hline
        $x$ & $-2$ & $-1{,}5$ & $-1$ & $-0{,}5$ & $0$ & $0{,}5$ & $1$ & $1{,}5$ & $2$\\
        \hline
        $f(x)$ & $4$ & $2{,}25$ & $1$ & $0{,}25$ & $0$ & $0{,}25$ & $1$ & $2{,}25$ & $4$\\
        \hline
        \end{tabular}}
  \item Courbe ci-dessous.
        \begin{center}
        \begin{tikzpicture}[x=0.6cm,y=0.6cm]
          \repere{-3}{-1}{3}{5}
          \gradx{-2,-1,1,2} \grady{1,2,3,4}
          \draw[coulCourbeB,thick,dashed] (-3,1) -- (3,1) node[above left,font=\scriptsize] {$y=1$};
          \draw[coulCourbeB,thick,dashed] (-3,4) -- (3,4) node[above left,font=\scriptsize] {$y=4$};
          \draw[coulCourbe,line width=1.1pt,domain=-2:2,samples=50] plot (\x,{\x*\x});
          \foreach \p in {(-2,4),(-1,1),(0,0),(1,1),(2,4)} \path[coulCourbe] \p node[croix]{};
          \node[coulCourbe,font=\footnotesize] at (1.4,3) {$\Cf$};
        \end{tikzpicture}
        \end{center}
  \item a) \res{S=\ens{-1\,;1}} \quad b) \res{S=\ens{-2\,;2}} \quad
        c) \res{S=\varnothing} : la courbe ne descend pas sous l'axe des abscisses.
  \item a) Courbe sur ou sous la droite $y=1$ : \res{S=\intff{-1}{1}}\\
        b) \res{S=\left[-2\,;-1\right[\cup\left]1\,;2\right]}\\
        c) \res{S=\left]-2\,;2\right[}
\end{questions}
\rem{On retrouve la propriété 6 : $x^{2}=1$ a deux solutions opposées. En 4) b), les
bornes $-2$ et $2$ sont comprises : ce sont les bornes de l'intervalle d'étude.}
\end{exo}

\begin{exo}[Se ramener à $x^{2}=k$]
\begin{questions}
  \item $5x^{2}+x^{2}=14+4$, soit $6x^{2}=18$ et $x^{2}=3$.
  \item $3>0$ : deux solutions, $-\sqrt{3}$ et $\sqrt{3}$.
  \item \res{S=\ens{-\sqrt{3}\,;\sqrt{3}}}
\end{questions}
\end{exo}

\begin{exo}[Équations avec des carrés]
\begin{reponses}
  \item $2x^{2}=18$, $x^{2}=9$ : \ \res{S=\ens{-3\,;3}}
  \item $-2x^{2}=-10$, $x^{2}=5$ : \ \res{S=\ens{-\sqrt{5}\,;\sqrt{5}}}
  \item $2x^{2}=-4$, $x^{2}=-2$ : \ \res{S=\varnothing}
\end{reponses}
\end{exo}

\partie{Fonction cube}

\begin{exo}[Images et antécédents par la fonction cube]
\begin{questions}
  \item \res{27} ; \ \res{-8} ; \ \res{0{,}008} ; \ \res{-125}
  \item \res{4} ; \ \res{-3} ; \ \res{10} ; \ \res{\sqrt[3]{6}}\,$\approx 1{,}82$
\end{questions}
\rem{Un seul antécédent à chaque fois, même pour un nombre négatif.}
\end{exo}

\begin{exo}[Lecture graphique avec la fonction cube]
\begin{questions}
  \item {\small\setlength{\tabcolsep}{3pt}\renewcommand{\arraystretch}{1.3}%
        \begin{tabular}{|c|c|c|c|c|c|c|c|}
        \hline
        $x$ & $-1{,}5$ & $-1$ & $-0{,}5$ & $0$ & $0{,}5$ & $1$ & $1{,}5$\\
        \hline
        $g(x)$ & $-3{,}4$ & $-1$ & $-0{,}1$ & $0$ & $0{,}1$ & $1$ & $3{,}4$\\
        \hline
        \end{tabular}}
  \item Courbe ci-dessous.
        \begin{center}
        \begin{tikzpicture}[x=0.7cm,y=0.7cm]
          \repere{-2}{-4}{2}{4}
          \gradx{-1,1} \grady{-3,-2,-1,1,2,3}
          \draw[coulCourbeB,thick,dashed] (-2,1) -- (2,1) node[above left,font=\scriptsize] {$y=1$};
          \draw[coulCourbeB,thick,dashed] (-2,-1) -- (2,-1) node[below left,font=\scriptsize] {$y=-1$};
          \draw[coulCourbe,line width=1.1pt,domain=-1.5:1.5,samples=50] plot (\x,{\x*\x*\x});
          \foreach \p in {(-1,-1),(0,0),(1,1)} \path[coulCourbe] \p node[croix]{};
          \node[coulCourbe,font=\footnotesize] at (1.75,3) {$\mathcal{C}_g$};
        \end{tikzpicture}
        \end{center}
  \item a) \res{S=\ens{1}} \quad b) \res{S=\ens{-1}} \quad c) \res{S=\ens{0}}
  \item a) \res{S=\left[-1{,}5\,;1\right]}\\
        b) \res{S=\left]-1\,;1{,}5\right]}\\
        c) \res{S=\left[0\,;1{,}5\right]}
\end{questions}
\rem{Une seule solution à chaque équation : la courbe monte tout le temps, chaque
droite horizontale la coupe une seule fois.}
\end{exo}

\begin{exo}[Équations avec des cubes (1)]
\begin{reponses}
  \item $-8x^{3}=-16$, $x^{3}=2$ : \ \res{S=\ens{\sqrt[3]{2}}}
  \item $3x^{3}=24$, $x^{3}=8$ : \ \res{S=\ens{2}}
\end{reponses}
\end{exo}

\begin{exo}[Équations avec des cubes (2)]
\begin{reponses}
  \item $3x^{3}=-24$, $x^{3}=-8$ : \ \res{S=\ens{-2}}
  \item $-3x^{3}=15$, $x^{3}=-5$ : \ \res{S=\ens{-\sqrt[3]{5}}}
\end{reponses}
\rem{Contrairement à $x^{2}=k$, l'équation $x^{3}=k$ a une solution même si $k<0$, et
une seule.}
\end{exo}

\partie{Fonction inverse}

\begin{exo}[Lecture graphique sur l'hyperbole]
\begin{questions}
  \item {\small\setlength{\tabcolsep}{3pt}\renewcommand{\arraystretch}{1.3}%
        \begin{tabular}{|c|c|c|c|c|c|c|c|c|}
        \hline
        $x$ & $-4$ & $-2$ & $-1$ & $-0{,}5$ & $0{,}5$ & $1$ & $2$ & $4$\\
        \hline
        $f(x)$ & $-0{,}25$ & $-0{,}5$ & $-1$ & $-2$ & $2$ & $1$ & $0{,}5$ & $0{,}25$\\
        \hline
        \end{tabular}}
  \item Courbe ci-dessous.
        \begin{center}
        \begin{tikzpicture}[x=0.6cm,y=0.6cm]
          \repere{-5}{-3}{5}{3}
          \gradx{-4,-3,-2,-1,1,2,3,4} \grady{-2,-1,1,2}
          \draw[coulCourbeB,thick,dashed] (-5,1) -- (5,1) node[above left,font=\scriptsize] {$y=1$};
          \draw[coulCourbeB,thick,dashed] (-5,-1) -- (5,-1) node[below left,font=\scriptsize] {$y=-1$};
          \draw[coulCourbe,line width=1.1pt,domain=0.5:4,samples=50] plot (\x,{1/\x});
          \draw[coulCourbe,line width=1.1pt,domain=-4:-0.5,samples=50] plot (\x,{1/\x});
          \foreach \p in {(-4,-0.25),(-1,-1),(-0.5,-2),(0.5,2),(1,1),(4,0.25)}
            \path[coulCourbe] \p node[croix]{};
          \node[coulCourbe,font=\footnotesize] at (1.3,2.3) {$\Cf$};
        \end{tikzpicture}
        \end{center}
  \item \res{S=\ens{1}} \quad et \quad \res{S=\ens{-1}}
  \item Courbe au-dessous de $y=1$ :
        \res{S=\left[-4\,;0\right[\cup\left]1\,;4\right]}\\
        Courbe sur ou au-dessus de $y=-1$ :
        \res{S=\left[-4\,;-1\right]\cup\left]0\,;4\right]}
\end{questions}
\rem{Pour $1/x<1$, toute la branche de gauche convient : pour $x<0$, $\frac{1}{x}$ est
négatif, donc inférieur à $1$. Le $0$ est toujours exclu : valeur interdite.}
\end{exo}

\begin{exo}[Équations et inéquations (1)]
\rem{On ne multiplie pas par $x$, dont le signe est inconnu : on se ramène à un
quotient comparé à $0$.}
\textbf{a)} $x\ne 0$ : $1=5x$ : \ \res{S=\ens{\tfrac{1}{5}}}\par
\textbf{b)} $\dfrac{1}{x}-5\le 0$, soit $\dfrac{1-5x}{x}\le 0$.
\begin{tabsignes}{$x$/1,$1-5x$/1,$x$/1,$\frac{1-5x}{x}$/1.2}{$-\infty$,$0$,$\frac{1}{5}$,$+\infty$}
  \tkzTabLine{,+,t,+,z,-,}
  \tkzTabLine{,-,z,+,t,+,}
  \tkzTabLine{,-,d,+,z,-,}
\end{tabsignes}
\res{S=\left]-\infty\,;0\right[\cup\left[\tfrac{1}{5}\,;+\infty\right[}\par
\textbf{c)} Même tableau, quotient $>0$ : \ \res{S=\left]0\,;\tfrac{1}{5}\right[}
\end{exo}

\begin{exo}[Équations et inéquations (2)]
\textbf{a)} $1=-2x$ : \ \res{S=\ens{-\tfrac{1}{2}}}\par
\textbf{b)} $\dfrac{1}{x}+2\le 0$, soit $\dfrac{1+2x}{x}\le 0$ ; numérateur nul en
$-\frac{1}{2}$, valeur interdite $0$ ; signe $+$, $-$, $+$ :\par
\res{S=\left[-\tfrac{1}{2}\,;0\right[}\par
\textbf{c)} \res{S=\left]-\tfrac{1}{2}\,;0\right[}
\end{exo}

\begin{exo}[Comparer deux inverses]
\begin{reponses}
  \item $\dfrac{1}{x}-\dfrac{1}{4}\le 0$, soit $\dfrac{4-x}{4x}\le 0$ ; signe $-$, $+$,
        $-$ (interdite $0$, numérateur nul en $4$) :\\
        \res{S=\left]-\infty\,;0\right[\cup\left[4\,;+\infty\right[}
  \item $\dfrac{5-3x}{5x}\ge 0$ ; signe $-$, $+$, $-$ (interdite $0$, numérateur nul
        en $\frac{5}{3}$) :\\
        \res{S=\left]0\,;\tfrac{5}{3}\right]}
  \item $\dfrac{2-x}{2x}>0$ : \ \res{S=\left]0\,;2\right[}
\end{reponses}
\rem{Erreur fréquente en a) : « $x\ge 4$ » en inversant les deux membres. On oublie
alors tous les $x$ négatifs.}
\end{exo}

\begin{exo}[Quotients du premier degré]
\begin{reponses}
  \item Interdite : $1$ ; $6=2(3x-3)$, soit $6=6x-6$ et $x=2$ : \ \res{S=\ens{2}}
  \item Interdite : $4$ ; $-4=2(-2x+8)$, soit $-4=-4x+16$ et $4x=20$ :
        \ \res{S=\ens{5}}
\end{reponses}
\rem{On vérifie que la solution n'est pas la valeur interdite.}
\end{exo}

\partie{Fonction racine carrée}

\begin{exo}[Lecture graphique avec la racine carrée (1)]
\begin{questions}
  \item {\small\setlength{\tabcolsep}{3.5pt}\renewcommand{\arraystretch}{1.3}%
        \begin{tabular}{|c|c|c|c|c|c|c|}
        \hline
        $x$ & $0$ & $0{,}5$ & $1$ & $2$ & $3$ & $4$\\
        \hline
        $f(x)$ & $0$ & $0{,}7$ & $1$ & $1{,}4$ & $1{,}7$ & $2$\\
        \hline
        \end{tabular}}
  \item Courbe ci-dessous.
        \begin{center}
        \begin{tikzpicture}[x=0.75cm,y=0.75cm]
          \repere{-1}{-1}{5}{3}
          \gradx{1,2,3,4} \grady{1,2}
          \draw[coulCourbeB,thick,dashed] (-1,1) -- (5,1) node[above left,font=\scriptsize] {$y=1$};
          \draw[coulCourbeB,thick,dashed] (-1,2) -- (5,2) node[above left,font=\scriptsize] {$y=2$};
          \draw[coulCourbe,line width=1.1pt,domain=0:4,samples=80] plot (\x,{sqrt(\x)});
          \foreach \p in {(0,0),(1,1),(4,2)} \path[coulCourbe] \p node[croix]{};
          \node[coulCourbe,font=\footnotesize] at (3,1.4) {$\Cf$};
        \end{tikzpicture}
        \end{center}
  \item a) \res{S=\ens{1}} \quad b) \res{S=\ens{4}} \quad c) \res{S=\ens{0}}
  \item a) \res{S=\intff{0}{1}} \quad b) \res{S=\left]1\,;4\right]} \quad
        c) \res{S=\left[0\,;4\right[}
\end{questions}
\rem{On retrouve la méthode 14 : $\sqrt{x}=2$ pour $x=2^{2}=4$. La courbe commence en
$0$ : jamais de $x$ négatif.}
\end{exo}

\begin{exo}[Lecture graphique avec la racine carrée (2)]
\begin{questions}
  \item {\small\setlength{\tabcolsep}{3pt}\renewcommand{\arraystretch}{1.3}%
        \begin{tabular}{|c|c|c|c|c|c|c|c|}
        \hline
        $x$ & $0$ & $0{,}5$ & $1$ & $2$ & $3$ & $4$ & $5$\\
        \hline
        $g(x)$ & $0$ & $1{,}4$ & $2$ & $2{,}8$ & $3{,}5$ & $4$ & $4{,}5$\\
        \hline
        \end{tabular}}
  \item Courbe ci-dessous.
        \begin{center}
        \begin{tikzpicture}[x=0.5cm,y=0.5cm]
          \repere{-1}{-1}{6}{6}
          \gradx{1,2,3,4,5} \grady{1,2,3,4,5}
          \draw[coulCourbeB,thick,dashed] (-1,2) -- (6,2) node[below left,font=\scriptsize] {$y=2$};
          \draw[coulCourbeB,thick,dashed] (-1,4) -- (6,4) node[below left,font=\scriptsize] {$y=4$};
          \draw[coulCourbeB,thick] (-1,-1) -- (6,6) node[below right,font=\scriptsize] {$y=x$};
          \draw[coulCourbe,line width=1.1pt,domain=0:5,samples=80] plot (\x,{2*sqrt(\x)});
          \foreach \p in {(0,0),(1,2),(4,4)} \path[coulCourbe] \p node[croix]{};
          \node[coulCourbe,font=\footnotesize] at (2,3.4) {$\mathcal{C}_g$};
        \end{tikzpicture}
        \end{center}
  \item \res{S=\ens{1}} \quad et \quad \res{S=\ens{4}}
  \item \res{S=\left]1\,;5\right]} \quad et \quad \res{S=\intff{0}{4}}
  \item La courbe et la droite se coupent en $(0\,;0)$ et en $(4\,;4)$ :
        \res{S=\ens{0\,;4}}. La courbe est sur ou au-dessus de la droite pour $x$
        entre $0$ et $4$ : \res{S=\intff{0}{4}}
\end{questions}
\end{exo}

\begin{exo}[Isoler la racine carrée (1)]
\rem{On isole $\sqrt{x}$, puis : $\sqrt{x}=k$ avec $k\ge 0$ donne $x=k^{2}$ ; avec
$k<0$, pas de solution.}
\begin{reponses}
  \item $2\sqrt{x}=8$, soit $\sqrt{x}=4$ : \ \res{S=\ens{16}}
  \item $3\sqrt{x}=9$, soit $\sqrt{x}=3$ : \ \res{S=\ens{9}}
  \item $2\sqrt{x}=-6$, soit $\sqrt{x}=-3$ : \ \res{S=\varnothing}
  \item $-3\sqrt{x}=-6$, soit $\sqrt{x}=2$ : \ \res{S=\ens{4}}
\end{reponses}
\rem{On vérifie en remplaçant : en a), $5\times 4+2=22$ et $3\times 4+10=22$.}
\end{exo}

\begin{exo}[Isoler la racine carrée (2)]
\begin{reponses}
  \item $4\sqrt{x}=20$, soit $\sqrt{x}=5$ : \ \res{S=\ens{25}}
  \item $3\sqrt{x}=-6$, soit $\sqrt{x}=-2$ : \ \res{S=\varnothing}
  \item $-3\sqrt{x}=-9$, soit $\sqrt{x}=3$ : \ \res{S=\ens{9}}
  \item $-3\sqrt{x}=0$, soit $\sqrt{x}=0$ : \ \res{S=\ens{0}}
\end{reponses}
\rem{En b), une racine carrée n'est jamais négative. En d), $\sqrt{x}=0$ a une
solution : $0$.}
\end{exo}

\begin{exo}[Inverses de carrés, de cubes et de racines]
\begin{reponses}
  \item Interdite : $0$ ; $4=x^{2}$ : \ \res{S=\ens{-2\,;2}}
  \item Interdite : $0$ ; $-2=16x^{3}$, $x^{3}=-\frac{1}{8}$ : \ \res{S=\ens{-\tfrac{1}{2}}}
  \item Il faut $x>0$ ; $21=3\sqrt{x}$, $\sqrt{x}=7$ : \ \res{S=\ens{49}}
\end{reponses}
\rem{En b), $\left(-\frac{1}{2}\right)^{3}=-\frac{1}{8}$ : on vérifie en remplaçant.}
\end{exo}

\end{document}
